\section{Data and the Anatomy of a Non-Market Transfer}\label{sec:data}

\subsection{Deeds}

Property records come from CoreLogic's Owner Transfer file
\citep{corelogic_data}, obtained through an academic license, which compiles
the universe of recorded deeds from county recorder offices. A deed is filed
whenever the legal ownership of a parcel changes---not only when a house is
sold. The same administrative trail that records a \$400,000 sale records a
mother quitclaiming her house to her son for nothing. Market-based datasets
(MLS listings, repeat-sales indices, mortgage originations) discard these
zero-price events by construction. Here they are the object of study.

The extract covers all fifty states and the District of Columbia. Coverage
is dense from 2007 through mid-2024; we use 2007--2023 for annual statistics
and treat 2024 as a partial year. Restricting to residential parcels and
deduplicating to one event per parcel per derived sale date yields 148.7 million
deed events (144.0 million in 2007--2023), linked through a permanent parcel
identifier that follows each house across successive owners.

\subsection{Classifying transfers}

CoreLogic flags each record with a primary category---arm's length or
not---and an interfamily indicator constructed from the recorder's
consideration, deed type, and name fields. We sharpen these flags into a
taxonomy using the buyer and seller names on the deed:

\begin{enumerate}
  \item \textbf{Market sale}: arm's-length category, recorded price of at
        least \$10,000, not flagged interfamily (71.1 million events,
        2007--2023).
  \item \textbf{Same-person retitle}: the named principal is the same person
        on both sides of the deed---an owner adding or removing a co-owner,
        potentially at marriage, divorce, or refinancing (7.1 million). The
        matched name remains, but ownership shares can change. These deeds
        are tracked separately from the operational family aggregate.
  \item \textbf{Family transfer, same surname}: flagged interfamily, buyer
        and seller sharing a surname but not the same person, no trust
        detected (13.6 million). This is a narrower proxy for inter-person
        family transfers: it misses unmatched surnames but can also retain
        same-person events with inconsistent full names. It is not a formal
        lower bound on true family transfers.
  \item \textbf{Family transfer, other}: flagged interfamily without a
        detected surname match after the earlier exclusions, plus the
        separately coded estate/executor residual (6.1 million). Missing
        surname fields and transfers involving entities can enter this group;
        the code does not verify the relationship between the parties.
  \item \textbf{Trust self-transfer (operational label)}: the interfamily
        flag is set and a primary party name contains a trust keyword
        (12.1 million). This can include estate-planning retitling and
        distributions to different beneficiaries. Unchanged control is not
        verified; the class is excluded from the family aggregate.
  \item \textbf{Other non-arm's-length}: the remainder---divorce divisions
        across surnames, title corrections, nominal-price conveyances, and
        foreclosure-related instruments (34.0 million).
\end{enumerate}

Table~\ref{tab:taxonomy} reports observable differences across the classes.
Same-surname family deeds differ sharply from the defined market-sales class: 90 percent of same-surname transfers record a zero or missing
price, against 0 percent of market sales (which require a price by
construction) and 65 percent of the residual non-arm's-length class; 40
percent travel by quitclaim deed---the instrument of gift---against 2 percent
of market sales. Where a price is recorded at all in a same-surname transfer,
the median is \$97,000, well under half the \$224,000 median market price:
the comparison is descriptive because the houses and transferred interests can differ. Corporate buyers, who take 10
percent of market purchases, are essentially absent from same-surname
transfers. These patterns are consistent with gifts and succession, but do not verify their prevalence or exclude other transactions. One
note on probate: formal estate keywords (``estate of,'' executor,
administrator) rarely appear in the recorded name fields, and the earlier trust, same-person, and surname rules also take precedence over the estate rule. The small residual estate class therefore cannot measure all successions.

\begin{table}[!t]\centering
\caption{Residential deed events by class, 2007--2023}
\label{tab:taxonomy}
{\footnotesize\setlength{\tabcolsep}{4pt}
\begin{tabular}{lrrrrr}
\toprule
 & N & Zero/no & Median & Quitclaim & Absentee \\
Event class & (2007--2023) & price (\%) & price$^{a}$ & share (\%) & recipient (\%)$^{b}$ \\
\midrule
Market sale (arm's length) & 71,076,255 & 0.0 & \$224k & 2.0 & 26.8 \\
Other non-arm's-length & 33,980,735 & 64.8 & \$57k & 17.8 & 55.9 \\
Family transfer: same surname & 13,610,459 & 90.3 & \$97k & 39.6 & 17.8 \\
Trust self-transfer & 12,074,757 & 96.1 & \$130k & 30.0 & 25.6 \\
Same-person retitle (co-owner change) & 7,123,579 & 92.3 & \$63k & 40.3 & 16.9 \\
Family transfer: other interfamily & 6,073,863 & 88.3 & \$155k & 26.4 & 34.8 \\
Estate deed (non-interfamily) & 57,697 & 9.6 & \$167k & 2.2 & 89.8 \\
Family transfer: estate/executor & 1,905 & 82.9 & \$596k & 8.6 & 35.6 \\
\bottomrule
\end{tabular}}
\begin{minipage}{0.95\textwidth}\vspace{2pt}{\scriptsize \textit{Notes:} Authors' classification of CoreLogic Owner Transfer records, residential parcels, deduplicated to one event per parcel-date. $^{a}$Median of strictly positive recorded prices. $^{b}$Share of events whose buyer mailing ZIP differs from the property ZIP, among events where both are observed.}\end{minipage}
\end{table}

